\section{Historical context from the source}
The following literature summary is retained from the supplied preprint;
it is not an independent survey of the current literature.

The conjecture was recorded by Dean and Latka~\cite{DeanLatka1995}.
A \emph{tournament} is an oriented graph with an arc between every pair of
distinct vertices.  This case, also known as Dean's conjecture, was proved
by Fisher~\cite{Fisher1996} using a probability distribution obtained from
Farkas' lemma.  Havet and Thomass\'e gave a combinatorial proof using
local median orders~\cite{HavetThomasse2000}.  Fidler and
Yuster~\cite{FidlerYuster2007} extended the tournament theorem to
tournaments missing a matching, a star, or a clique.  Ghazal corrected the
star-case proof and treated generalized stars~\cite{Ghazal2012}.
Subsequent advances include a matching together with a
star~\cite{DaraFrancisJacobNarayanan2022} and two
stars~\cite{DaamouchAlMninyGhazal2025}.

Other approaches impose degree conditions or weaken the desired ratio.
Kaneko and Locke established the conjecture for minimum outdegree at most
six~\cite{KanekoLocke2001}; a recent computer-assisted preprint treats
minimum outdegree seven~\cite{SadhukhanSandeepSen2026}.
Chen, Shen and Yuster proved a universal second-to-first neighborhood
ratio of approximately $0.657298$~\cite{ChenShenYuster2003}.
Huang and Peng raise this to approximately $0.715538$ and additionally
report a computational bound of $0.74530$~\cite{HuangPeng2024}.
Botler, Moura and Naia~\cite{BotlerMouraNaia2023}, followed by
Espuny D\'iaz, Gir\~ao, Granet and Kronenberg~\cite{EspunyDiazEtAl2025},
studied every orientation of binomial random graphs; the latter result
covers every fixed edge probability below $1/2$, with probability tending
to one as the order tends to infinity.  Brukhman's recent dense-case
preprint concerns graphs of order $n\leq 2\delta+2$, where $\delta$ is
the minimum outdegree~\cite{Brukhman2026}.
